\documentclass{article}
\usepackage[UTF8]{ctex}
\usepackage{amsmath,amssymb,amsfonts}
\usepackage{enumitem}
\usepackage{geometry}
\geometry{a4paper, margin=1in}
% === 错题本排版优化（不参与题目标识） ===
\setlength{\parindent}{0pt}
\setlength{\parskip}{0.35em plus 0.10em minus 0.08em}
\setlength{\emergencystretch}{3em}
\setlist{itemsep=0.25em, topsep=0.40em, parsep=0pt}
\allowdisplaybreaks[2]
% === 错题本排版优化结束 ===
\begin{document}
\section*{华中科技大学2025年数学分析真题}
\section*{133．华中科技大学2025年数学分析试卷}

\subsection{题目1}
1．求极限 \(\lim _{x \rightarrow 0} \frac{\left(1+x^{2}\right)^{3}-\cos ^{2} x}{\tan x^{2}}\) ．

\subsection{题目2}
2．讨论级数 \(\sum_{n=1}^{\infty} \int_{0}^{\frac{\pi}{n}} \frac{\sin x}{1+x^{2024}} \mathrm{~d} x\) 的敛散性．

\subsection{题目3}
3．设 \(\beta>0, \alpha>1\) 为常数，计算无穷积分 \(I=\int_{0}^{+\infty} \frac{\mathrm{d} x}{1+\beta x^{\alpha}}\) ．

\subsection{题目4}
4．设 \(\alpha>0\) ，求极限 \(\lim _{x \rightarrow+\infty}\left[(x+1)^{\alpha}-x^{\alpha}-\frac{x+1}{x+e^{\theta \ln x}}\right]\) ．

\subsection{题目5}
5．求极限 \(\lim _{n \rightarrow \infty} \sum_{k=1}^{n} \frac{n^{2}}{n^{3}+k^{3}+k^{2}}\) ．

\subsection{题目6}
6．求极限 \(\lim _{x \rightarrow 0^{+}} \frac{e^{x}-1-x}{\sqrt{1-x}-\cos \sqrt{x}}\) ．

\subsection{题目7}
7．求积分

\[
I=\oint_{C}(x+y)^{2} \mathrm{~d} x-\left(x^{2}+y^{2}\right) \mathrm{d} y .
\]

其中 \(C\) 是从点 \(A(1,1)\) 到点 \(B(2,5)\) 再到点 \(C(3,2)\) 最后到点 \(A(1,1)\) 的三角形边界．

\subsection{题目8}
8．设 \(a, b, c\) 为正常数，\(S\) 为单位球面 \(x^{2}+y^{2}+z^{2}=1\) 的外侧，计算第二类曲面积分

\[
I=\iint_{S} \frac{x \mathrm{~d} y \mathrm{~d} z+y \mathrm{~d} z \mathrm{~d} x+z \mathrm{~d} x \mathrm{~d} y}{\left(a x^{2}+b y^{2}+c z^{2}\right)^{\frac{3}{2}}} .
\]

\subsection{题目9}
9．用 \(\varepsilon-N\) 语言证明极限 \(\lim _{n \rightarrow \infty} \frac{1}{\sqrt[n]{(2 n)!}}=0\) ．

\subsection{题目10}
10．设

\[
\begin{gathered}
x_{1}=1, x_{2}=\sqrt{\frac{1}{2}+1}, x_{3}=\sqrt{\frac{1}{3}+\sqrt{\frac{1}{2}+1}}, x_{4}=\sqrt{\frac{1}{4}+\sqrt{\frac{1}{3}+\sqrt{\frac{1}{2}+1}}}, \\
\cdots, x_{n+1}=\sqrt{\frac{1}{n+1}+x_{n}}, \cdots
\end{gathered}
\]

证明数列 \(\left\{x_{n}\right\}\) 收敛并求出极限值。

\subsection{题目11}
11．设函数 \(f(x) \in C(0,+\infty)\) ，满足 \(x=f(x) 5^{f(x)}\) ．证明：

\begin{itemize}
  \item[（1）] \(f(x)\) 单调递增．
  \item[（2）] \(\lim _{x \rightarrow+\infty} f(x)=+\infty\) ．
  \item[（3）] \(\lim _{x \rightarrow+\infty} \frac{f(x)}{\ln x}=\frac{1}{\ln 5}\) ．
\end{itemize}

\subsection{题目12}
12．对常数 \(a>0\) ，记平面区域 \(D: x^{2}+y^{2} \leq a^{2}\) 的边界为 \(\partial D\) 。设二元函数 \(f(x, y)\) 在 \(D\) 上有连续偏导数，且 \(f(x, y)=0,(x, y) \in \partial D\) ．证明：

\[
\left|\iint_{D} f(x, y) \mathrm{d} x \mathrm{~d} y\right| \leq \frac{\pi a^{2}}{3} \max _{(x, y) \in D} \sqrt{\left(\frac{\partial f}{\partial x}\right)^{2}+\left(\frac{\partial f}{\partial y}\right)^{2}} .
\]
\end{document}
